% 开发日志章节
\section{开发日志}

本章用于记录项目开发过程中的关键事件、技术方案、思考过程以及待办事项，为团队协作和后续论文撰写提供详细的历史资料。每个里程碑对应一个日志条目，条目下按时间或任务划分子条目。

\subsection{里程碑一：调研与需求分析日志}

本日志记录了第一个里程碑阶段的活动，重点在于梳理需求、研究现有系统并形成规范文档。

\subsubsection{任务分解与进展}

\begin{itemize}
  \item \textbf{文献与源码研读：}团队成员阅读了 ABIDES 等现有市场模拟器的论文和开源代码，梳理其核心设计理念。重点关注 ABIDES 的消息驱动设计、纳秒级时间精度、可配置网络延迟等特性。
  \item \textbf{需求收集：}通过对接科研人员和交易员，收集了跨资产仿真、动态校准、可视化交互等需求点。对需求进行了优先级划分，并记录在需求规格说明书中。
  \item \textbf{竞品分析：}对 MAXE、JAX-LOB 等模拟器进行功能对比和性能评估，分析其在并行度、数据驱动特性方面的优势与不足。
  \item \textbf{需求文档撰写：}整理上述调研内容，形成初版需求文档，包括功能需求、非功能需求、约束条件和潜在风险。
\end{itemize}

\subsubsection{技术路径与方案选择}

\begin{itemize}
  \item \textbf{架构选择：}根据需求与调研结论，决定采用事件驱动的分层架构，将系统划分为核心仿真层、服务支撑层与交互层，并采用消息总线实现组件解耦。
  \item \textbf{语言与并发模型：}综合考虑开发效率与性能，初步选定 Python 为主要语言，内核采用多线程或异步结合的模式，以提升事件处理吞吐量。
  \item \textbf{数据结构：}针对订单簿实现，研究了跳表、红黑树等数据结构的复杂度，初步倾向于跳表以便于实现并发友好的有序队列。
  \item \textbf{校准算法：}调研了基于误差反馈的 PID 控制、自适应元胞等自校准方案，为后续设计提供思路。
\end{itemize}

\subsubsection{问题与思考}

\begin{itemize}
  \item 在对 ABIDES 的解析中发现，其核心内核单线程架构对多资产模拟存在性能瓶颈。我们需要在保持时间同步的基础上设计适合并行扩展的调度器。
  \item 需求收集过程中，利益相关者期望系统能同时支撑策略测试和政策仿真，这对数据接口和交互层的灵活性提出了挑战。
  \item 如何在早期需求尚不稳定时制定可行的设计路线，并在后续迭代中保持架构的可扩展性，是需要持续思考的问题。
\end{itemize}

\subsubsection{代办与改进}

\begin{itemize}
  \item 深入研究高性能数据结构库，评估第三方并行订单簿实现的可行性。
  \item 与利益相关者进一步讨论需求细节，明确跨资产策略和自校准指标的具体定义。
  \item 构建最小可运行原型（MVP）的计划，验证核心内核与订单簿设计的可行性。
\end{itemize}

\subsubsection{总结}

在此阶段，我们完成了对现有模拟器的全面调研和需求分析，形成了完整的需求说明书和初步的技术路线。通过深入理解 ABIDES 等系统的优点与不足，我们确定了采用事件驱动、分层架构和并行化内核的方向，并针对订单簿数据结构、自校准算法等关键技术进行了初步选型。后续阶段将基于这些结论开展系统设计和原型实现。